\section{Proof of Theorem \ref{thm form}}\label{sec PROOF}

\subsection{Formula (\ref{id B}) follows from formula (\ref{id A})}
The Lie algebra $\widehat{{\mathfrak{sl}_2}}$ has an automorphism $\rho$, corresponding to the involution of the Dynkin diagram:
\begin{gather*}
\rho(e_{i})=e_{3-i}, \qquad \rho(f_{i})=f_{3-i}, \qquad \rho(h_{i})=h_{3-i}, \qquad i=1,2.
\end{gather*}
We have $\rho^2={\rm id}$. In other words, $\rho$ acts by the formulas
\begin{gather*}
e\leftrightarrow fT, \qquad f \leftrightarrow e T^{-1}, \qquad h \leftrightarrow c-h.
\end{gather*}

\begin{lem}\label{lem elts}For $i \in \mathbb{Z}_{> 0}$, we have
\begin{gather*}
\rho \colon \ \frac{f}{T^i} \mapsto \frac{e}{T^{i+1}}, \qquad \frac{e}{T^i} \mapsto \frac{f}{T^{i-1}}, \qquad \frac{h}{T^i} \mapsto -\frac{h}{T^i}.
\end{gather*}
\end{lem}
\begin{proof}
We have
\begin{gather*}
\frac{f}{T} = \frac12 \left[f, \frac hT \right] = \frac 12 \left[ f, \left[\frac eT, f \right] \right] \overset{\delta}{\longrightarrow} \frac 12 \left[ \frac eT, \left[f, \frac eT \right] \right] = \frac 12 \left[ \frac eT, - \frac hT \right] = \frac{e}{T^2}, \nonumber \\
\frac{f}{T^i} = \frac 12 \left[ \frac{f}{T^{i-1}}, \left[\frac{e}{T}, f\right] \right] \overset{\delta}{\longrightarrow} \frac 12 \left[ \frac{e}{T^i}, \left[f, \frac eT \right] \right] = \frac 12 \left[ \frac{e}{T^i}, - \frac hT \right] = \frac{e}{T^{i+1}}. \nonumber
\end{gather*}
Similarly we prove that $\rho\big(\frac{e}{T^i}\big) = \frac{f}{T^{i-1}}$, $\rho \big(\frac{h}{T^i}\big) = -\frac{h}{T^i}$.
\end{proof}

For $m\in{\mathbb C}$, let $\sigma_m \colon \widehat{{\mathfrak{sl}_2}} \to \operatorname{End}( V(m,k-m))$ be the Verma module structure. Let $\sigma_m\circ \rho \colon \widehat{{\mathfrak{sl}_2}} \to \operatorname{End}( V(m,k-m))$ be the twisted module structure.

Clearly the $\widehat{{\mathfrak{sl}_2}}$-modules $(\sigma_m\circ \rho, V(m,k-m))$ and $(\sigma_{m-k}, V(m-k,m))$ are isomorphic. If $v_m\in V(m,k-m)$ and $v_{k-m}\in V(k-m,m)$ are generating vectors, then an isomorphism $\chi\colon (\sigma_m\circ \rho, V(m,k-m)) \to (\sigma_{m-k}, V(m-k,m))$ is defined by the formula,
\begin{gather*}
f_{i_l}\cdots f_{i_1}v_{k-m} \mapsto f_{3-i_l}\cdots f_{3-i_1}v_{m},
\end{gather*}
for any $i_1,\dots,i_l\in\{1,2\}$. The isomorphism $\chi$ restricts to isomorphisms of the graded components, $V(k-m,m)_{(p_1,p_2)} \to V(m,k-m)_{(p_2,p_1)}$.

In Section \ref{BASIS} we fixed bases of the homogeneous components $V_{(p_1,p_2)}$ with $p_1\ne p_2$ of any Verma module $V$. By Lemma~\ref{lem elts}, under the isomorphism $\chi$ the chosen basis of $V(k-m, m)_{(p_1,p_2)}$ is mapped to the chosen basis of $V(m,k-m)_{(p_2,p_1)}$ up to multiplication of the basis vectors by~$\pm 1$. This~$\pm 1$ appears due to the formula $\rho\big(\frac h{T^i}\big) = - \frac h{T^i}$. In particular, we have
\begin{gather*}
\chi \colon \ \frac f{T^i}v_{k-m} \mapsto \frac e{T^{i+1}}v_{m}, \qquad
\frac f{T^i} \frac f{T^j} v_{k-m} \mapsto \frac e{T^{i+1}}\frac e{T^{j+1}}v_{m}.
\end{gather*}

Let $ \sigma^*_m \colon \widehat{{\mathfrak{sl}_2}} \to \operatorname{End}( V(m,k-m)^*)$ be the contragradient Verma module structure. Let $ \sigma^*_m\circ \rho \colon \widehat{{\mathfrak{sl}_2}} \to \operatorname{End}( V(m,k-m)^*)$ be the twisted module structure. The isomorphism $\chi$ induces an isomorphism of modules $\chi^* \colon (\sigma^*_m\circ \rho, V(m,k-m)^*) \to (\sigma^*_{m-k}, V(m-k,m)^*)$.

In Section~\ref{BASIS} we fixed bases in the homogeneous components $V^*_{(p_1,p_2)}$ with $p_1\ne p_2$ of any contragradient Verma module~$V^*$. Under the isomorphism $\chi^*$, the chosen basis of $V(k-m, m)^*_{(p_1,p_2)}$ is mapped to the chosen basis of $V(m,k-m)^*_{(p_2,p_1)}$ up to multiplication of the basis vectors by~$\pm 1$. In particular, we have
\begin{gather*}
 \chi^* \colon \ \left(\frac f{T^i}v_{k-m}\right)^* \mapsto \left(\frac e{T^{i+1}}v_{m}\right)^*, \qquad
\left(\frac f{T^i} \frac f{T^j} v_{k-m}\right)^* \mapsto \left(\frac e{T^{i+1}}\frac e{T^{j+1}}v_{m}\right)^*.
\end{gather*}

Assume that the relation in formula (\ref{id A}) holds in every contragradient Verma module $V^*$. Then in $V(k-m, m)^*$ it takes the form
\begin{gather*}
\frac{f}{T^{a-1}} (v_{k-m})^* = (-m-2+ a(k+2))\left(\frac{f}{T^{a-1}}v_{k-m}\right)^* \nonumber \\
\hphantom{\frac{f}{T^{a-1}} (v_{k-m})^* =}{} + \sum_{\ell=1}^{a-1} \bigg[\frac{h}{T^\ell}\left(\frac{f}{T^{a-1-\ell}}v_{k-m}\right)^* + 2 \frac{e}{T^\ell}\mathop{\sum_{i+j=a-1-\ell}}_{i\geq j\geq 0}
\left(\frac{f}{T^i}\frac{f}{T^j}v_{k-m}\right)^*\bigg].%\label{id Am}
\end{gather*}
The isomorphism $\chi^*$ sends this relation to the relation in $V(m,k-m)^*$,
\begin{gather*}
\frac{e}{T^{a}} (v_{m})^* = (-m-2+ a(k+2))\left(\frac{e}{T^{a}}v_{m}\right)^* \nonumber\\
\hphantom{\frac{e}{T^{a}} (v_{m})^* =}{} + \sum_{\ell=1}^{a-1} \bigg[{-}\frac{h}{T^\ell}\left(\frac{e}{T^{a-\ell}}v_{m}\right)^* + 2 \frac{f}{T^{\ell-1}}\mathop{\sum_{i+j=a-1-\ell}}_{i\geq j\geq 0}
\left(\frac{e}{T^{i+1}}\frac{e}{T^{j+1}}v_{m}\right)^*\bigg],%\label{id Amm}
\end{gather*}
which is exactly the relation in formula~(\ref{id B}). Thus formula (\ref{id A}) implies formula~(\ref{id B}).

\subsection{Auxiliary lemma}\label{subsec aux lemma}
Let
\begin{gather*}
V=V(m,k-m) \qquad \text{and} \qquad V^* = V(m,k-m)^*.
\end{gather*}

\begin{lem} \label{lem pere} For $x\in V$, $\phi\in V^*$, $k\in{\mathbb Z}_{\geq 0}$, we have
\begin{gather*}
\left\langle \frac{f}{T^k} \varphi, x \right\rangle = \big\langle \varphi, eT^k x \big\rangle,
 \qquad \left\langle \frac{e}{T^k} \varphi, x \right\rangle = \big\langle \varphi, fT^k x \big\rangle, \qquad
\left\langle \frac{h}{T^k} \varphi, x \right\rangle = \big\langle \varphi, hT^k x \big\rangle. %\label{perek}
\end{gather*}
\end{lem}
\begin{proof}The proof is by induction. We prove the first equality, the others are proved similarly.

We have $[f_2, f_1] = \frac{h}{T}$, hence $[f_1, [f_2, f_1]] = \frac{2f}{T}$. Similarly
 $[e_1, [e_2, e_1]] = 2e T$. So for $k=1$, we have
\begin{gather*}
\left\langle \frac{f}{T} \varphi, x \right\rangle = \left\langle \frac{1}{2} [f_1, [f_2, f_1]] \varphi, x \right\rangle = \left\langle \varphi, \frac{1}{2} [[e_1, e_2], e_1] x \right\rangle = \langle \varphi, eT x \rangle.
\end{gather*}
We have $\big[f_2, \frac{f}{T^{k-1}}\big] = \frac{h}{T^k}$, hence $\big[f_1, \big[f_2, \frac{f}{T^{k-1}} \big]\big] = \frac{2f}{T^k}$. Similarly, $\big[\big[eT^{k-1},fT\big], e\big] = 2e T^k$. Then
\begin{gather*}
\left\langle \frac{f}{T^k} \varphi, x \right\rangle = \left\langle \frac 12 \left[f_1, \left[f_2, \frac{f}{T^{k-1}}\right]\right] \varphi, x \right\rangle = \big\langle \varphi, \big[\big[eT^{k-1},e_2\big], e_1\big] x \big\rangle = \big\langle \varphi, eT^k x \big\rangle.\tag*{\qed}
\end{gather*}\renewcommand{\qed}{}
\end{proof}


\subsection{The structure of the proof of formula (\ref{id A})}\label{subsec structure} We reformulate formula (\ref{id A}) as
\begin{gather}
(m+(a-1)(k+2))\left(\frac{f}{T^{a-1}}v\right)^* \nonumber\\
\qquad{} =\frac{f}{T^{a-1}} (v)^* - \sum_{\ell=1}^{a-1}\bigg[\frac{h}{T^\ell}\left(\frac{f}{T^{a-1-\ell}}v\right)^* + 2 \frac{e}{T^\ell}\mathop{\sum_{i+j=a-1-\ell}}_{i\geq j\geq 0} \left(\frac{f}{T^i}\frac{f}{T^j}v\right)^*\bigg], \label{AA}
\end{gather}
and will prove it in this form.

Each term in (\ref{AA}) is an element of the homogeneous component $V^*_{(a,a-1)}$. In Section~\ref{BASIS} we specified a basis of the dual component $V_{(a,a-1)}$. We will calculate the value of the right-hand side in~(\ref{AA}) on an arbitrary basis vector and will obtain the value of the left-hand side on that vector.

The basis in $V_{(a,a-1)}$ consists of the vectors
\begin{gather*}
\frac{f}{T^{i_1}}\cdots\frac{f}{T^{i_r}} \frac{h}{T^{j_1}}\cdots\frac{h}{T^{j_s}} \frac{e}{T^{l_1}}\cdots\frac{e}{T^{l_{r-1}}}v, %\label{bb}
\end{gather*}
where
\begin{gather*}
0\leq i_r\leq i_{r-1}\leq\dots\leq i_1, \qquad 1\leq j_s\leq j_{s-1}\leq\dots\leq j_1,\qquad 1\leq l_{r-1}\leq l_{r-2}\leq\dots\leq l_1;\nonumber\\
 \sum_{u=1}^r i_u + \sum_{u=1}^s j_u + \sum_{u=1}^{r-1} l_u = a-1. %\label{indice}
\end{gather*}
We partition the basis in four groups. Group~${\rm O}$ consists of the single basis vector $\frac{f}{T^{a-1}}v$. Group~${\rm I}$ consists of all basis vectors with $r=1$, but different from $\frac{f}{T^{a-1}}v$. Group ${\rm II}$ consists of all basis vectors with $r=2$. Group ${\rm III}$ consists of all basis vectors with $r\geq 3$.

Notice that the value of the left-hand side of~(\ref{AA}) on the basis vector~$\frac{f}{T^{a-1}}v$ equals $m + (a-1)(k+2)$. Hence we need to show that the value of the right-hand side on the basis vector $\frac{f}{T^{a-1}}v$ equals $m + (a-1)(k+2)$. Similarly the value of the left-hand side on any basis vector of Groups~\mbox{I--III} equals zero. Hence we need
to prove that the value of the right-hand side on any basis vector of Groups~I--III equals zero. These four statements are the content of Propositions~\ref{Pr4},~\ref{Pr3},~\ref{Pr2}, and~\ref{Pr1} below. These propositions prove Theorem~\ref{thm form}.

\subsection[Group ${\rm O}$]{Group $\boldsymbol{{\rm O}}$}\label{subsec group O}
\begin{prop}\label{Pr4}The value of the right-hand side of \eqref{AA} on the basis vector $\frac{f}{T^{a-1}}v$ equals $m + (a-1)(k+2)$.
\end{prop}
\begin{proof}By Lemma \ref{lem pere} we have{\samepage
\begin{gather*}
\left\langle \frac{f}{T^{a-1}} (v)^*, \frac{f}{T^{a-1}} v \right\rangle = \left\langle (v)^*, eT^{a-1} \frac{f}{T^{a-1}} v \right\rangle \\
\hphantom{\left\langle \frac{f}{T^{a-1}} (v)^*, \frac{f}{T^{a-1}} v \right\rangle} = \left\langle (v)^*, \left[ h+(a-1)c + \frac{f}{T^{a-1}} eT^{a-1} \right] v \right\rangle = m + (a-1)k,
\end{gather*}
since $eT^{a-1} v$ is of degree $(-a, -a+1)$, hence zero.}

By Lemma \ref{lem pere}, for $ \ell \in \{ 1, \dots, a-1 \}$ we have
\begin{gather*}
 \bigg\langle \frac{h}{T^{\ell}} \left( \frac{f}{T^{a-1-\ell}} v \right)^*, \frac{f}{T^{a-1}} v \bigg\rangle
 = \bigg\langle \left( \frac{f}{T^{a-1-\ell}} v \right)^*, hT^{\ell} \frac{f}{T^{a-1}} v \bigg\rangle \\
 \hphantom{\bigg\langle \frac{h}{T^{\ell}} \left( \frac{f}{T^{a-1-\ell}} v \right)^*, \frac{f}{T^{a-1}} v \bigg\rangle}{} = \bigg\langle \left( \frac{f}{T^{a-1-\ell}} v \right)^*, -2 \frac{f}{T^{a-1-\ell}} v \bigg\rangle = -2.
\end{gather*}
By Lemma \ref{lem pere} for $ \ell \in \{ 1, \dots, a-1 \}$ we have
\begin{gather*}
\bigg\langle \frac{e}{T^{\ell}} \mathop{\sum_{i+j=a-1-\ell}}_{i\geq j\geq 0} \left(\frac{f}{T^i} \frac{f}{T^j}v \right)^*, \frac{f}{T^{a-1}} v \bigg\rangle = \bigg\langle \mathop{\sum_{i+j=a-1-\ell}}_{i\geq j\geq 0} \left(\frac{f}{T^i} \frac{f}{T^j}v \right)^*, fT^{\ell} \frac{f}{T^{a-1}} v \bigg\rangle \\
\hphantom{\bigg\langle \frac{e}{T^{\ell}} \mathop{\sum_{i+j=a-1-\ell}}_{i\geq j\geq 0} \left(\frac{f}{T^i} \frac{f}{T^j}v \right)^*, \frac{f}{T^{a-1}} v \bigg\rangle} = \bigg\langle \mathop{\sum_{i+j=a-1-\ell}}_{i\geq j\geq 0} \left(\frac{f}{T^i} \frac{f}{T^j}v \right)^*, \frac{f}{T^{a-1}} fT^{\ell} v \bigg\rangle =0,
\end{gather*}
since $fT^{\ell} v$ is of degree $ (-\ell+1, -\ell) \leq (0,-1)$, hence zero. Therefore,
\begin{gather*}
\bigg\langle \frac{f}{T^{a-1}} (v)^* - \sum_{\ell=1}^{a-1} \bigg[\frac{h}{T^\ell}\left(\frac{f}{T^{a-1-\ell}}v\right)^* + 2\frac{e}{T^\ell} \mathop{\sum_{i+j=a-1-\ell}}_{i\geq j\geq 0} \left(\frac{f}{T^i}\frac{f}{T^j}v\right)^* \bigg], \frac{f}{T^{a-1}} v \bigg\rangle \\
\qquad{}= m + (a-1)k - (-2)(a-1) = m + (a-1)(k+2).
\end{gather*}
Proposition \ref{Pr4} is proved.
\end{proof}

\subsection{Group I}
\begin{prop}\label{Pr3}The value of the right-hand side of \eqref{AA} on any basis vector of Group ${\rm I}$ equals zero.
\end{prop}
\begin{proof}Group ${\rm I}$ consists of basis vectors of the form
\begin{gather*}
w = \frac{f}{T^{a-1-n}}\frac{h}{T^{j_1}} \cdots \frac{h}{T^{j_s}} v, \qquad \text{where} \quad
n \in \{1, \dots, a-1 \},\quad j_1 + \dots + j_s = n, \quad j_i \geq 1.
\end{gather*}

\begin{lem}\label{Group 3, s=1} In the notation above, if $s =1$, then
\begin{gather}
\bigg\langle \frac{f}{T^{a-1}} (v)^*, w \bigg\rangle = 2 n k, \nonumber\\ %\label{G1s11}\\
\bigg\langle \frac{h}{T^{\ell}} \left(\frac{f}{T^{a-1-\ell}}v\right)^*, w \bigg\rangle = \begin{cases}
 2 n k, & \text{if $\ell=n$,}\\
 - 4, & \text{if $\ell > a-1-n$,} \\
 0, & \text{if $\ell \leq a-1-n$,}
 \end{cases} \label{G1s12}
\\
\bigg\langle \frac{e}{T^{\ell}} \mathop{\sum_{i+j=a-1-\ell}}_{i\geq j\geq 0} \left(\frac{f}{T^i}\frac{f}{T^j}v\right)^*, w \bigg\rangle =
 \begin{cases}
 2, & \text{if $\ell \leq n$,}\\
 0, & \text{if $\ell > n$. }
 \end{cases} \nonumber %\label{G1s13}
\end{gather}
\end{lem}

Note that the first line in (\ref{G1s12}) is not mutually exclusive with the second and third lines in~(\ref{G1s12}).

\begin{proof}
We have $w = \frac{f}{T^{a-1-n}}\frac{h}{T^n} v $. Then
\begin{gather*}
\bigg\langle \frac{f}{T^{a-1}} (v)^*, \frac{f}{T^{a-1-n}}\frac{h}{T^n} v \bigg\rangle = \bigg\langle (v)^*, eT^{a-1} \frac{f}{T^{a-1-n}}\frac{h}{T^n} v \bigg\rangle \\
\hphantom{\bigg\langle \frac{f}{T^{a-1}} (v)^*, \frac{f}{T^{a-1-n}}\frac{h}{T^n} v \bigg\rangle }
= \bigg\langle (v)^*, \left[ hT^n + \frac{f}{T^{a-1-n}} eT^{a-1} \right] \frac{h}{T^n} v \bigg\rangle.
\end{gather*}
Note that $eT^{a-1} \frac{h}{T^n}v$ is of degree $(-a+n, -a+1+n) \leq (-1,0)$, so that $eT^{a-1} \frac{h}{T^n} v = 0$. Hence
\begin{gather*}
\bigg\langle (v)^*, hT^n \frac{h}{T^n} v \bigg\rangle = \bigg\langle (v)^*, \left[2nk + \frac{h}{T^n} hT^n\right] v \bigg\rangle = 2 n k.
\end{gather*}
We have
\begin{gather*}
\bigg\langle \frac{h}{T^{\ell}} \left(\frac{f}{T^{a-1-\ell}}v\right)^*,
\frac{f}{T^{a-1-n}}\frac{h}{T^n} v \bigg\rangle = \bigg\langle \left(\frac{f}{T^{a-1-\ell}}v\right)^*, hT^{\ell} \frac{f}{T^{a-1-n}}\frac{h}{T^n} v \bigg\rangle,\\
hT^{\ell} \frac{f}{T^{a-1-n}}\frac{h}{T^n} v = \left[ -2 fT^{\ell+n-a+1} + \frac{f}{T^{a-1-n}} hT^{\ell} \right] \frac{h}{T^n} v \\
\hphantom{hT^{\ell} \frac{f}{T^{a-1-n}}\frac{h}{T^n} v} = -2 fT^{\ell+n-a+1} \frac{h}{T^n} v + \frac{f}{T^{a-1-n}} hT^{\ell} \frac{h}{T^n} v.
\end{gather*}
Note that the second summand is nonzero if and only if $\ell=n$. In that case we have
\begin{gather*}
\bigg\langle \left(\frac{f}{T^{a-1-n}}v\right)^*, \frac{f}{T^{a-1-n}} hT^n \frac{h}{T^n} v \bigg\rangle = 2nk.
\end{gather*}
For the first summand, if $\ell+n-a+1 \leq 0$, then $-2 fT^{\ell+n-a+1} \frac{h}{T^n} v$ is a basis vector and so pairing with $\left(\frac{f}{T^{a-1-\ell}}v\right)^*$ gives zero. If $\ell+n-a+1 > 0$, then
\begin{gather*}
\bigg\langle \left(\frac{f}{T^{a-1-\ell}}v\right)^*, -2 fT^{\ell+n-a+1} \frac{h}{T^n} v \bigg\rangle \\
\qquad {} = \bigg\langle \left(\frac{f}{T^{a-1-\ell}}v\right)^*, -2 \left[ 2\frac{f}{T^{a-1-\ell}} + \frac{h}{T^n} fT^{\ell+n-a+1}\right] v \bigg\rangle = - 4,
\end{gather*}
where we used $fT^{\ell+n-a+1}v = 0$.

Finally,
\begin{gather*}
\bigg\langle \frac{e}{T^{\ell}} \mathop{\sum_{i+j=a-1-\ell}}_{i\geq j\geq 0} \left(\frac{f}{T^i}\frac{f}{T^j}v\right)^*, \frac{f}{T^{a-1-n}}\frac{h}{T^n} v \bigg\rangle = \bigg\langle \mathop{\sum_{i+j=a-1-\ell}}_{i\geq j\geq 0} \left(\frac{f}{T^i}\frac{f}{T^j}v\right)^*, fT^{\ell} \frac{f}{T^{a-1-n}}\frac{h}{T^n} v \bigg\rangle, \\
fT^{\ell} \frac{f}{T^{a-1-n}}\frac{h}{T^n} v = \frac{f}{T^{a-1-n}} fT^{\ell} \frac{h}{T^n} v = \frac{f}{T^{a-1-n}} \left[ 2fT^{\ell-n} + \frac{h}{T^n} fT^{\ell} \right] v = 2 \frac{f}{T^{a-1-n}} \frac{f}{T^{n-\ell}} v,
\end{gather*}
since $ fT^{\ell} v = 0$. Note that $(a-1-n)+(n-\ell)=a-1-\ell$, hence if $i = a-1-n$ and $j = n-\ell$ (or vice versa depending on what is greater) we have
\begin{gather*}
\bigg\langle \mathop{\sum_{i+j=a-1-\ell}}_{i\geq j\geq 0} \left(\frac{f}{T^i}\frac{f}{T^j}v\right)^*, 2 \frac{f}{T^{a-1-n}} \frac{f}{T^{n-\ell}} v \bigg\rangle = 2,
\end{gather*}
whenever $n-\ell \geq 0$ and zero otherwise. The lemma is proved.
\end{proof}

For $s=1$ Proposition \ref{Pr3} follows from Lemma~\ref{Group 3, s=1}:
\begin{gather*}
\bigg\langle \frac{f}{T^{a-1}} (v)^* - \sum_{\ell=1}^{a-1} \bigg[\frac{h}{T^{\ell}}\left(\frac{f}{T^{a-1-\ell}}v\right)^* + 2\frac{e}{T^{\ell}} \mathop{\sum_{i+j=a-1-\ell}}_{i\geq j\geq 0} \left(\frac{f}{T^i}\frac{f}{T^j}v\right)^*\bigg], \frac{f}{T^{a-1-n}}\frac{h}{T^n} v \bigg\rangle \\
\qquad = 2nk - 2nk + 4n - 2 \cdot 2n = 0.
\end{gather*}

\begin{lem}\label{lem G1s}
For $s \geq 2$, we have
\begin{gather}
\bigg\langle \frac{f}{T^{a-1}}(v) ^*, w \bigg\rangle =0, \label{G1sg1} \\
\sum_{\ell=1}^{a-1} \bigg\langle \frac{h}{T^{\ell}} \left(\frac{f}{T^{a-1-\ell}}v\right)^*, w \bigg\rangle =0, \label{G1sg2} \\
\sum_{\ell=1}^{a-1} \bigg\langle \frac{e}{T^{\ell}} \mathop{\sum_{i+j=a-1-\ell}}_{i\geq j\geq 0} \left(\frac{f}{T^i}\frac{f}{T^j}v\right)^*, w \bigg\rangle = 0. \label{G1sg3}
\end{gather}
\end{lem}

\begin{proof}Recall that $w=\frac{f}{T^{a-1-n}}\frac{h}{T^{j_1}} \cdots \frac{h}{T^{j_s}} v$ with $j_1+\cdots+j_s=n$. We have
\begin{gather*}
\bigg\langle \frac{f}{T^{a-1}} (v)^*, w \bigg\rangle = \big\langle (v)^*, eT^{a-1} w \big\rangle, \\
eT^{a-1} w = eT^{a-1} \frac{f}{T^{a-1-n}}\frac{h}{T^{j_1}} \cdots \frac{h}{T^{j_s}} v = \left[ hT^n + \frac{f}{T^{a-1-n}} eT^{a-1} \right] \frac{h}{T^{j_1}} \cdots \frac{h}{T^{j_s}} v.
\end{gather*}
We have $hT^n \frac{h}{T^{j_1}} \cdots \frac{h}{T^{j_s}} v=0$, since $hT^n$ commutes with all $\frac{h}{T^{j_i}}$. Indeed, we have $n > j_i$ since $j_1 + \dots + j_s = n$, $j_i \geq 1$, and $s \geq 2$.

We also have $\frac{f}{T^{a-1-n}} eT^{a-1} \frac{h}{T^{j_1}} \cdots \frac{h}{T^{j_s}} v=0$ since $eT^{a-1} \frac{h}{T^{j_1}} \cdots \frac{h}{T^{j_s}}v$ is of degree $(-a+n, -a+1+n) \leq (-1, 0)$, hence zero. This proves~(\ref{G1sg1}).

We prove (\ref{G1sg2}) by induction on $s$. For $s=2$ we have
\begin{gather*}
\bigg\langle \frac{h}{T^{\ell}} \left(\frac{f}{T^{a-1-\ell}}v\right)^*, w \bigg\rangle = \bigg\langle \left(\frac{f}{T^{a-1-\ell}}v\right)^*, hT^{\ell} w \bigg\rangle, \\
hT^{\ell} w = hT^{\ell} \frac{f}{T^{a-1-n}}\frac{h}{T^{j_1}} \frac{h}{T^{j_2}} v = \left[ -2 fT^{\ell-a+1+n} + \frac{f}{T^{a-1-n}} hT^{\ell} \right] \frac{h}{T^{j_1}} \frac{h}{T^{j_2}} v \\
\hphantom{hT^{\ell} w}{} = -2 fT^{\ell-a+1+n} \frac{h}{T^{j_1}} \frac{h}{T^{j_2}} v + \frac{f}{T^{a-1-n}} hT^{\ell} \frac{h}{T^{j_1}} \frac{h}{T^{j_2}} v.
\end{gather*}
Note that for $ \frac{f}{T^{a-1-n}} hT^{\ell} \frac{h}{T^{j_1}} \frac{h}{T^{j_2}} v $ to give a nonzero pairing with $\big(\frac{f}{T^{a-1-\ell}}v\big)^*$ we need $\ell=n$, which implies that $hT^{\ell}$ commutes with $\frac{h}{T^{j_1}}$ and $\frac{h}{T^{j_2}}$ ($\ell > j_i$ since $j_1+j_2 =n=\ell$ and $j_i \geq 1$), so that $ \frac{f}{T^{a-1-n}} hT^{\ell} \frac{h}{T^{j_1}} \frac{h}{T^{j_2}} v $ gives zero for all $\ell$.

Also note that whenever $\ell \leq a-1-n$, $fT^{\ell-a+1+n} \frac{h}{T^{j_1}} \frac{h}{T^{j_2}} v$ is a basis vector and so pairing with $\big(\frac{f}{T^{a-1-\ell}}v\big)^*$ gives zero. If $\ell > a-1-n$, then
\begin{gather}
-2 fT^{\ell-a+1+n} \frac{h}{T^{j_1}} \frac{h}{T^{j_2}} v = -2 \left[ 2 fT^{\ell-a+1+n-j_1} + \frac{h}{T^{j_1}} fT^{\ell-a+1+n} \right] \frac{h}{T^{j_2}} v \nonumber \\
\hphantom{-2 fT^{\ell-a+1+n} \frac{h}{T^{j_1}} \frac{h}{T^{j_2}} v}{} = -4 fT^{\ell-a+1+n-j_1} \frac{h}{T^{j_2}} v -2 \frac{h}{T^{j_1}} fT^{\ell-a+1+n} \frac{h}{T^{j_2}} v. \label{1 and 2}
\end{gather}

If $\ell \leq a-1-n+j_1$, the first summand gives zero when pairing with $\big(\frac{f}{T^{a-1-\ell}}v\big)^*$, since for such $\ell fT^{\ell-a+1+n-j_1} \frac{h}{T^{j_2}} v$ is a basis vector. For $\ell > a-1-n+j_1$, we have
\begin{gather*}
-4 fT^{\ell-a+1+n-j_1} \frac{h}{T^{j_2}} v = -4 \left[ 2 \frac{f}{T^{a-1-\ell}} v + \frac{h}{T^{j_2}} fT^{\ell-a+1+n-j_1} \right] v= -8 \frac{f}{T^{a-1-\ell}} v,
\end{gather*}
since $fT^{\ell-a+1+n-j_1}v$ is of degree $(-\ell+a-1-n + j_1 +1, -\ell+a-1-n+j_1) \leq (0, -1)$, hence must be equal to zero. So for $\ell \in \{a-1-n+j_1 +1, \dots, a-1\}$ we get
\begin{gather*}
\bigg\langle \left(\frac{f}{T^{a-1-\ell}}v\right)^*, -4 fT^{\ell-a+1+n-j_1} \frac{h}{T^{j_2}} v \bigg\rangle = -8
\end{gather*}
and zero for other values of $\ell$. The total number of elements in the set $\{a-1-n+j_1 +1, \dots, a-1\}$ equals $j_2$.

Whereas, for the second summand in (\ref{1 and 2}) we have
\begin{gather*}
-2 \frac{h}{T^{j_1}} fT^{\ell-a+1+n} \frac{h}{T^{j_2}} v = -2\frac{h}{T^{j_1}} \left[ 2 fT^{\ell-a+1+n - j_2} + \frac{h}{T^{j_2}} fT^{\ell-a+1+n} \right] v \\
\hphantom{-2 \frac{h}{T^{j_1}} fT^{\ell-a+1+n} \frac{h}{T^{j_2}} v } = -4 \frac{h}{T^{j_1}} fT^{\ell-a+1+n -j_2} v,
\end{gather*}
since $fT^{\ell-a+1+n} v$ is of degree $(-\ell+a-1-n+1, -\ell+a-1-n) \leq (0, -1)$, hence must be equal to zero.

If $\ell > a-1-n+j_2$, then $fT^{\ell-a+1+n-j_2} v$ is of degree $(-\ell+a-1-n+j_2 +1, -\ell+a-1-n+j_2) \leq (0, -1)$, hence must be equal to zero. If $\ell \leq a-1-n+j_2$, then
\begin{gather*}
-4 \frac{h}{T^{j_1}} fT^{\ell-a+1+n -j_2} v = -4 \left[ -2\frac{f}{T^{a-1-\ell}} + fT^{\ell-a+1+n - j_2} \frac{h}{T^{j_1}} \right] \\
\hphantom{-4 \frac{h}{T^{j_1}} fT^{\ell-a+1+n -j_2} v} = 8 \frac{f}{T^{a-1-\ell}} -4 fT^{\ell-a+1+n - j_2} \frac{h}{T^{j_1}}.
\end{gather*}
The second summand gives zero when pairing with $\big(\frac{f}{T^{a-1-\ell}}v\big)^*$. So for $\ell \in \{a-1-n+1, \dots, a-1-n+j_2 \}$ we get
\begin{gather*}
\bigg\langle \left(\frac{f}{T^{a-1-\ell}}v\right)^*, -2 \frac{h}{T^{j_1}} fT^{\ell-a+1+n} \frac{h}{T^{j_2}} v \bigg\rangle = 8
\end{gather*}
and zero for other values of $\ell$. The total number of elements in the set $\{a-1-n+1, \dots, a-1-n+j_2 \}$ equals $j_2$.

Therefore,
\begin{gather*}
\sum_{\ell=1}^{a-1} \bigg\langle \frac{h}{T^{\ell}} \left(\frac{f}{T^{a-1-\ell}}v\right)^*, \frac{f}{T^{a-1-n}} \frac{h}{T^{j_1}} \frac{h}{T^{j_2}} v \bigg\rangle = -8 j_2 + 8 j_2 = 0
\end{gather*}
and so for $s=2$ we proved (\ref{G1sg2}).

Now suppose that (\ref{G1sg2}) holds for all natural numbers up to $s$. Then
\begin{gather*}
hT^{\ell} \frac{f}{T^{a-1-n}} \frac{h}{T^{j_1}}\dots\frac{h}{T^{j_{s+1}}} v = \left[ -2 fT^{\ell-a+1+n} + \frac{f}{T^{a-1-n}} hT^{\ell} \right] \frac{h}{T^{j_1}}\cdots\frac{h}{T^{j_{s+1}}} v.
\end{gather*}
Note that for $\frac{f}{T^{a-1-n}} hT^{\ell} \frac{h}{T^{j_1}} \cdots \frac{h}{T^{j_{s+1}}} v$ to give a nonzero pairing with $\big(\frac{f}{T^{a-1-\ell}}v\big)^*$
we need $\ell=n$. That assumption implies that $hT^{\ell}$ commutes with $\frac{h}{T^{j_i}}$ for all $i \in \{1, \dots, s+1\}$ since $\ell > j_i$ as $j_1+\dots +j_{s+1} =n=\ell$ and $j_i \geq 1$. Hence $ \frac{f}{T^{a-1-n}} hT^{\ell} \frac{h}{T^{j_1}}\cdots \frac{h}{T^{j_{s+1}}} v $ gives zero for all $\ell$.

Also note that whenever $\ell \leq a-1-n$, the vector $fT^{\ell-a+1+n} \frac{h}{T^{j_1}} \cdots \frac{h}{T^{j_{s+1}}} v$ is a basis vector and so pairing with $\big(\frac{f}{T^{a-1-\ell}}v\big)^*$ gives zero.

If $\ell > a-1-n$, then
\begin{gather}
-2 fT^{\ell-a+1+n} \frac{h}{T^{j_1}}\cdots\frac{h}{T^{j_{s+1}}} v=-2 \left[ 2fT^{\ell-a+1+n - j_1} + \frac{h}{T^{j_1}} fT^{\ell-a+1+n}\right] \frac{h}{T^{j_2}}\cdots\frac{h}{T^{j_{s+1}}} v \nonumber \\
\qquad{} = -4 \frac{f}{T^{a-1 - (n-j_1) - \ell}}\frac{h}{T^{j_2}}\cdots\frac{h}{T^{j_{s+1}}}v - 2 \frac{h}{T^{j_1}} fT^{-a+1 + n + \ell} \frac{h}{T^{j_2}}\cdots\frac{h}{T^{j_{s+1}}}v. \label{fs sum}
\end{gather}

Note that by induction hypothesis we have
\begin{gather*}
0 = \sum_{\ell=1}^{a-1} \bigg\langle \frac{h}{T^{\ell}} \left(\frac{f}{T^{a-1-\ell}}v\right)^*, \frac{f}{T^{a-1-(n-j_1)}} \frac{h}{T^{j_2}}\cdots\frac{h}{T^{j_{s+1}}} v \bigg\rangle \nonumber \\
\hphantom{0} = \sum_{\ell=1}^{a-1} \bigg\langle \left(\frac{f}{T^{a-1-\ell}}v\right)^*, hT^{\ell} \frac{f}{T^{a-1-(n-j_1)}} \frac{h}{T^{j_2}}\cdots\frac{h}{T^{j_{s+1}}} v \bigg\rangle. %\label{temporary}
\end{gather*}

So we add this zero term multiplied by $-2$ to the first summand in $(\ref{fs sum})$ to get
\begin{gather}
 -2 \sum_{\ell=1}^{a-1} \bigg\langle \left(\frac{f}{T^{a-1-\ell}}v\right)^*, 2\frac{f}{T^{a-1 - (n-j_1) -\ell}}\frac{h}{T^{j_2}}\cdots\frac{h}{T^{j_{s+1}}} v \bigg\rangle \nonumber \\
\qquad\quad{} -2 \sum_{\ell=1}^{a-1} \bigg\langle \left(\frac{f}{T^{a-1-\ell}}v\right)^*, hT^{\ell} \frac{f}{T^{a-1-(n-j_1)}} \frac{h}{T^{j_2}}\cdots\frac{h}{T^{j_{s+1}}} v \bigg\rangle \label{trick_notes} \\
\qquad{} = -2 \sum_{\ell=1}^{a-1} \bigg\langle \left(\frac{f}{T^{a-1-\ell}}v\right)^*, \left[ 2\frac{f}{T^{a-1 - (n-j_1) - \ell}} + hT^{\ell} \frac{f}{T^{a-1-(n-j_1)}} \right] \frac{h}{T^{j_2}}\dots\frac{h}{T^{j_{s+1}}} v \bigg\rangle \nonumber \\
\qquad{} = -2 \sum_{\ell=1}^{a-1} \bigg\langle \left(\frac{f}{T^{a-1-\ell}}v\right)^*, \frac{f}{T^{a-1-(n-j_1)}} hT^{\ell} \frac{h}{T^{j_2}}\cdots\frac{h}{T^{j_{s+1}}} v \bigg\rangle, \nonumber
\end{gather}
where in the last step we use commutation relations.

Note that for $ \frac{f}{T^{a-1-(n-j_1)}} hT^{\ell} \frac{h}{T^{j_2}} \cdots \frac{h}{T^{j_{s+1}}} v $ to give a nonzero pairing with $\big(\frac{f}{T^{a-1-\ell}}v\big)^*$ we need $\ell=n-j_1$. That assumption implies that $hT^{\ell}$ commutes with $\frac{h}{T^{j_i}}$ for all $i \in \{2, \dots, s+1\}$ since $\ell>j_i$ as $j_2+\dots+ j_{s+1} =n - j_1=\ell$ and $j_i \geq 1$. Hence $\frac{f}{T^{a-1-(n-j_1)}} hT^{\ell} \frac{h}{T^{j_2}} \cdots \frac{h}{T^{j_{s+1}}} v $ gives zero for all~$\ell$.

For the second summand in (\ref{fs sum}) we have
\begin{gather}
- 2 \sum_{\ell=1}^{a-1} \bigg\langle \left(\frac{f}{T^{a-1-\ell}}v\right)^*, \frac{h}{T^{j_1}} fT^{-a+1 + n + \ell} \frac{h}{T^{j_2}}\cdots\frac{h}{T^{j_{s+1}}}v \bigg\rangle
\nonumber \\
\qquad{} = - 2 \sum_{\ell=1}^{a-1} \bigg\langle \left(\frac{f}{T^{a-1-\ell}}v\right)^*, \frac{h}{T^{j_1}} \left[ 2 \frac{f}{T^{a-1 - (n-j_2) - \ell}} + \frac{h}{T^{j_2}} fT^{-a+1 + n +\ell} \right] \frac{h}{T^{j_3}} \cdots\frac{h}{T^{j_{s+1}}}v \bigg\rangle \nonumber \\
\qquad{} = - 4 \sum_{\ell=1}^{a-1} \bigg\langle \left(\frac{f}{T^{a-1-\ell}}v\right)^*, \frac{h}{T^{j_1}} \frac{f}{T^{a-1 - (n-j_2) - \ell}} \frac{h}{T^{j_3}} \cdots\frac{h}{T^{j_{s+1}}}v \bigg\rangle \label{ind1} \\
\qquad\quad{} - 2 \sum_{\ell=1}^{a-1} \bigg\langle \left(\frac{f}{T^{a-1-\ell}}v\right)^*, \frac{h}{T^{j_1}} \frac{h}{T^{j_2}} fT^{-a+1 + n +\ell} \frac{h}{T^{j_3}} \cdots\frac{h}{T^{j_{s+1}}}v \bigg\rangle. \label{ind2}
\end{gather}

In (\ref{ind1}) we note that
\begin{gather*}
\frac{h}{T^{j_1}} \frac{f}{T^{a-1 - (n-j_2) - \ell}} \frac{h}{T^{j_3}} \cdots\frac{h}{T^{j_{s+1}}}v \\
\qquad{} = \left[ -2 \frac{f}{T^{a-1 - (n-j_1-j_2) - \ell}} + \frac{f}{T^{a-1 - (n-j_2) - \ell}}\frac{h}{T^{j_1}} \right] \frac{h}{T^{j_3}} \cdots\frac{h}{T^{j_{s+1}}}v \\
\qquad{} = -2 \frac{f}{T^{a-1 - (n-j_1-j_2) - \ell}} \frac{h}{T^{j_3}} \cdots\frac{h}{T^{j_{s+1}}}v + \frac{f}{T^{a-1 - (n-j_2) - \ell}}\frac{h}{T^{j_1}} \frac{h}{T^{j_3}} \cdots\frac{h}{T^{j_{s+1}}}v.
\end{gather*}
Note that in both terms the number of $h$'s is less than or equal to $s$, so we use the exact same reasoning as in (\ref{trick_notes}) to show that
\begin{gather*}
\sum_{\ell=1}^{a-1} \bigg\langle \left(\frac{f}{T^{a-1-\ell}}v\right)^*, \frac{f}{T^{a-1 - (n-j_1-j_2) - \ell}} \frac{h}{T^{j_3}} \dots\frac{h}{T^{j_{s+1}}}v \bigg\rangle = 0, \\
\sum_{\ell=1}^{a-1} \bigg\langle \left(\frac{f}{T^{a-1-\ell}}v\right)^*, \frac{f}{T^{a-1 - (n-j_2) - \ell}}\frac{h}{T^{j_1}} \frac{h}{T^{j_3}} \cdots\frac{h}{T^{j_{s+1}}}v \bigg\rangle = 0,
\end{gather*}
which implies that the expression in (\ref{ind1}) equals zero. Similarly, one shows that in (\ref{ind2}),
\begin{gather*}
\sum_{\ell=1}^{a-1} \bigg\langle \left(\frac{f}{T^{a-1-\ell}}v\right)^*, \frac{h}{T^{j_1}} \frac{h}{T^{j_2}} fT^{-a+1 + n + \ell} \frac{h}{T^{j_3}} \cdots\frac{h}{T^{j_{s+1}}}v \bigg\rangle
\end{gather*}
the factor $fT^{-a+1 + n + \ell}$ can be pulled to the right by using the same argument (first commute $fT^{-a+1 + n + \ell}$ with $\frac{h}{T^{j_3}}$ and then pull $fT^{-a+1 + n-j_3 + \ell}$ to the left). Ultimately, we get
\begin{gather*}
\sum_{\ell=1}^{a-1} \bigg\langle \left(\frac{f}{T^{a-1-\ell}}v\right)^*, \frac{h}{T^{j_1}} \frac{h}{T^{j_2}} \frac{h}{T^{j_3}} \cdots\frac{h}{T^{j_{s+1}}} fT^{-a+1 + n + \ell} v \bigg\rangle = 0,
\end{gather*}
since $\ell > a-1-n$ and so $fT^{-a+1 + n + \ell} v = 0$. Therefore,
\begin{gather*}
\sum_{\ell=1}^{a-1} \bigg\langle \frac{h}{T^{\ell}} \left(\frac{f}{T^{a-1-\ell}}v\right)^*, \frac{f}{T^{a-1-n}} \frac{h}{T^{j_1}}\cdots\frac{h}{T^{j_{s}}} v \bigg\rangle = 0,
\end{gather*}
and formula (\ref{G1sg2}) is proved.

We prove formula (\ref{G1sg3}) by induction on $s$. For $s=2$, we have
\begin{gather*}
\bigg\langle \frac{e}{T^{\ell}} \mathop{\sum_{i+j=a-1-\ell}}_{i\geq j\geq 0} \left(\frac{f}{T^i}\frac{f}{T^j}v\right)^*, w \bigg\rangle = \bigg\langle \mathop{\sum_{i+j=a-1-\ell}}_{i\geq j\geq 0} \left(\frac{f}{T^i}\frac{f}{T^j}v\right)^*, fT^{\ell} w \bigg\rangle, \\
 fT^{\ell} w = fT^{\ell} \frac{f}{T^{a-1-n}}\frac{h}{T^{j_1}} \frac{h}{T^{j_2}} v = \frac{f}{T^{a-1-n}} fT^{\ell} \frac{h}{T^{j_1}} \frac{h}{T^{j_2}} v.
\end{gather*}
Note that $fT^{\ell} \frac{h}{T^{j_1}} \frac{h}{T^{j_2}} v$ is of degree $(n-\ell+1, n-\ell)$, hence nonzero only if $\ell \leq n$. For such $\ell$ we have
\begin{gather}
\frac{f}{T^{a-1-n}} \left[ 2fT^{\ell-j_1} + \frac{h}{T^{j_1}} ft^{\ell} \right] \frac{h}{T^{j_2}} v = 2 \frac{f}{T^{a-1-n}} fT^{\ell-j_1} \frac{h}{T^{j_2}} v + \frac{f}{T^{a-1-n}} \frac{h}{T^{j_1}} fT^{\ell} \frac{h}{T^{j_2}} v \nonumber \\
\hphantom{\frac{f}{T^{a-1-n}} \left[ 2fT^{\ell-j_1} + \frac{h}{T^{j_1}} ft^{\ell} \right] \frac{h}{T^{j_2}} v} = 2 \frac{f}{T^{a-1-n}} fT^{\ell-j_1} \frac{h}{T^{j_2}} v +2 \frac{f}{T^{a-1-n}} \frac{h}{T^{j_1}} fT^{\ell-j_2} v. \label{2terms}
\end{gather}
If $\ell \leq j_1$, then the first summand in~(\ref{2terms}) gives zero when pairing with any vector with two $f$'s. If $\ell > j_1$, then
\begin{gather*}
2 \frac{f}{T^{a-1-n}} fT^{\ell-j_1} \frac{h}{T^{j_2}} v = 2 \frac{f}{T^{a-1-n}} \left[ 2fT^{\ell-j_1-j_2} + \frac{h}{T^{j_2}}fT^{\ell-j_1} \right] v = 4 \frac{f}{T^{a-1-n}} \frac{f}{T^{n-\ell}} v.
\end{gather*}
If $\ell > j_2$, then the second summand in (\ref{2terms}) is zero simply because $fT^{\ell-j_2} v = 0$. If $\ell \leq j_2$, then
\begin{gather*}
2 \frac{f}{T^{a-1-n}} \frac{h}{T^{j_1}} fT^{\ell-j_2} v = 2 \frac{f}{T^{a-1-n}} \left[ -2 fT^{\ell-j_1-j_2} + fT^{\ell-j_2}\frac{h}{T^{j_1}} \right] v = -4 \frac{f}{T^{a-1-n}} \frac{f}{T^{n-\ell}} v,
\end{gather*}
since $\frac{f}{T^{a-1-n}} fT^{\ell-j_2}\frac{h}{T^{j_1}}v$ is a basis vector, hence pairing with a vector consisting of two $f$'s gives zero. Therefore,
\begin{gather}
\sum_{\ell=1}^{a-1} \bigg\langle \mathop{\sum_{i+j=a-1-\ell}}_{i\geq j\geq 0} \left(\frac{f}{T^i}\frac{f}{T^j}v\right)^*, fT^{\ell} w \bigg\rangle \nonumber \\
\qquad{} = \sum_{\ell=1}^{a-1} \bigg\langle \mathop{\sum_{i+j=a-1-\ell}}_{i\geq j\geq 0}
\left(\frac{f}{T^i}\frac{f}{T^j}v\right)^*, 4 \frac{f}{T^{a-1-n}} \frac{f}{T^{n-\ell}} v \bigg\rangle \label{sum1} \\
\qquad\quad{} + \sum_{\ell=1}^{a-1} \bigg\langle \mathop{\sum_{i+j=a-1-\ell}}_{i\geq j\geq 0} \left(\frac{f}{T^i}\frac{f}{T^j}v\right)^*, -4 \frac{f}{T^{a-1-n}} \frac{f}{T^{n-\ell}} v \bigg\rangle. \label{sum2}
\end{gather}
Note that in the expression in~(\ref{sum1}) for each $\ell \in \{j_1 +1, \dots, n \}$ there exists exactly one pair of indices $(i,j) = (\max \{a-1-n, n-\ell \} , \min \{a-1-n, n-\ell\})$ that gives $4$ when pairing. All other pairs $(i,j)$ give zero. Similarly, the expression in~(\ref{sum2}) equals $-4$ for each $\ell \in \{1, \dots, j_2 \}$ and exactly one corresponding pair $(i,j)$, and zero otherwise. Also note that the number of elements in each set $\{j_1 +1, \dots, n \}$ and $\{1, \dots, j_2 \}$ equals $j_2$. Hence we get
\begin{gather*}
4 j_2 - 4 j_2 = 0.
\end{gather*}
Therefore, formula (\ref{G1sg3}) is proved for $s=2$.

Now suppose that formula (\ref{G1sg3}) holds for all natural numbers up to $s$. Then
\begin{gather}
fT^{\ell} \frac{f}{T^{a-1-n}} \frac{h}{T^{j_1}} \cdots \frac{h}{T^{j_{s+1}}} v\ =\
 \frac{f}{T^{a-1-n}} fT^{\ell} \frac{h}{T^{j_1}} \cdots \frac{h}{T^{j_{s+1}}} v \nonumber \\
\qquad{} = \frac{f}{T^{a-1-n}} \left[ 2 fT^{\ell-j_1} + \frac{h}{T^{j_1}} fT^{\ell} \right] \frac{h}{T^{j_2}} \cdots \frac{h}{T^{j_{s+1}}} v \nonumber \\
\qquad{} = 2 \frac{f}{T^{a-1-n}} fT^{\ell-j_1} \frac{h}{T^{j_2}} \cdots \frac{h}{T^{j_{s+1}}} v + \frac{f}{T^{a-1-n}} \frac{h}{T^{j_1}} fT^{\ell} \frac{h}{T^{j_2}} \cdots \frac{h}{T^{j_{s+1}}} v. \label{fs sum2}
\end{gather}

Note that if $\ell \leq j_1$, then the first summand in (\ref{fs sum2}) is a basis vector and hence its pairing with a vector consisting of two $f$'s gives zero. If $\ell > j_1$ we have
\begin{gather*}
\sum_{\ell=1}^{a-1} \bigg\langle \mathop{\sum_{i+j=a-1-\ell}}_{i\geq j\geq 0} \left(\frac{f}{T^i}\frac{f}{T^j}v\right)^*, \frac{f}{T^{a-1-n}} fT^{\ell-j_1} \frac{h}{T^{j_2}} \cdots \frac{h}{T^{j_{s+1}}} v \bigg\rangle\\
\qquad{} = \sum_{\ell=j_1+1}^{a-1} \bigg\langle \mathop{\sum_{i+j=a-1-\ell}}_{i\geq j\geq 0} \left(\frac{f}{T^i}\frac{f}{T^j}v\right)^*, \frac{f}{T^{a-1-n}} fT^{\ell-j_1} \frac{h}{T^{j_2}} \cdots \frac{h}{T^{j_{s+1}}} v \bigg\rangle \\
\qquad{} = \sum_{\ell=j_1+1}^{a-1} \bigg\langle \frac{e}{T^{\ell-j_1}} \mathop{\sum_{i+j=a-1-\ell}}_{i\geq j\geq 0} \left(\frac{f}{T^i}\frac{f}{T^j}v\right)^*, \frac{f}{T^{a-1-n}} \frac{h}{T^{j_2}} \cdots \frac{h}{T^{j_{s+1}}} v \bigg\rangle \\
\qquad{} = \sum_{k=1}^{a-1-j_1} \bigg\langle \frac{e}{T^{k}} \mathop{\sum_{i+j=a-1-j_1 - k}}_{i\geq j\geq 0} \left(\frac{f}{T^i}\frac{f}{T^j}v\right)^*, \frac{f}{T^{(a-1-j_1)-(n-j_1)}} \frac{h}{T^{j_2}} \cdots \frac{h}{T^{j_{s+1}}} v \bigg\rangle = 0
\end{gather*}
by induction hypothesis. For the second summand in (\ref{fs sum2}) we have
\begin{gather}
\frac{f}{T^{a-1-n}} \frac{h}{T^{j_1}} fT^{\ell} \frac{h}{T^{j_2}} \cdots \frac{h}{T^{j_{s+1}}} v = \frac{f}{T^{a-1-n}} \frac{h}{T^{j_1}} \left[ 2fT^{\ell-j_2} + \frac{h}{T^{j_2}} fT^{\ell} \right] \frac{h}{T^{j_3}} \cdots \frac{h}{T^{j_{s+1}}} v \nonumber \\
\qquad{} = 2 \frac{f}{T^{a-1-n}} \frac{h}{T^{j_1}} fT^{\ell-j_2} \frac{h}{T^{j_3}} \cdots \frac{h}{T^{j_{s+1}}} v + \frac{f}{T^{a-1-n}} \frac{h}{T^{j_1}} \frac{h}{T^{j_2}} fT^{\ell} \frac{h}{T^{j_3}} \cdots \frac{h}{T^{j_{s+1}}} v. \label{2f}
\end{gather}
Note that
\begin{gather*}
\frac{f}{T^{a-1-n}} \frac{h}{T^{j_1}} fT^{\ell-j_2} \frac{h}{T^{j_3}} \cdots \frac{h}{T^{j_{s+1}}} v = \frac{f}{T^{a-1-n}} \left[ -2 fT^{\ell-j_1-j_2} + fT^{\ell-j_2} \frac{h}{T^{j_1}} \right] \frac{h}{T^{j_3}} \cdots \frac{h}{T^{j_{s+1}}} v \\
\qquad{} = -2 \frac{f}{T^{a-1-n}} fT^{\ell-j_1-j_2} \frac{h}{T^{j_3}} \cdots \frac{h}{T^{j_{s+1}}} v + \frac{f}{T^{a-1-n}} fT^{\ell-j_2} \frac{h}{T^{j_1}} \frac{h}{T^{j_3}} \cdots \frac{h}{T^{j_{s+1}}} v,
\end{gather*}
where in each vector the number of $h$'s is less than or equal to~$s$. Repeating the argument above, we see that by induction hypothesis we get
\begin{gather*}
\sum_{\ell=1}^{a-1} \bigg\langle \mathop{\sum_{i+j=a-1-\ell}}_{i\geq j\geq 0} \left(\frac{f}{T^i}\frac{f}{T^j}v\right)^*, \frac{f}{T^{a-1-n}} \frac{h}{T^{j_1}} fT^{\ell-j_2} \frac{h}{T^{j_3}} \cdots \frac{h}{T^{j_{s+1}}} v \bigg\rangle = 0.
\end{gather*}
Now in the second summand in (\ref{2f}),
\begin{gather*}
\frac{f}{T^{a-1-n}} \frac{h}{T^{j_1}} \frac{h}{T^{j_2}} fT^{\ell} \frac{h}{T^{j_3}} \cdots \frac{h}{T^{j_{s+1}}} v,
\end{gather*}
we pull $fT^{\ell}$ to the right and at each step we use induction hypothesis to argue that we keep getting zeros. Ultimately, we get a vector
\begin{gather*}
\frac{f}{T^{a-1-n}} \frac{h}{T^{j_1}} \cdots \frac{h}{T^{j_{s+1}}} fT^{\ell} v,
\end{gather*}
which is zero, since $fT^{\ell}$ has grading $(-\ell+1, -\ell) \leq (0,-1)$ and so $fT^{\ell}v=0$. Therefore,
\begin{gather*}
\sum_{\ell=1}^{a-1} \bigg\langle \frac{e}{T^{\ell}} \mathop{\sum_{i+j=a-1-\ell}}_{i\geq j\geq 0} \left(\frac{f}{T^i}\frac{f}{T^j}v\right)^*, \frac{f}{T^{a-1-n}} \frac{h}{T^{j_1}} \cdots \frac{h}{T^{j_{s}}} v \bigg\rangle = 0.
\end{gather*}
Formula (\ref{G1sg3}) and Lemma \ref{lem G1s} are proved.
\end{proof}
Proposition \ref{Pr3} is proved.
\end{proof}

\subsection[Group ${\rm II}$]{Group $\boldsymbol{{\rm II}}$}\label{subsec group 2}
\begin{prop}\label{Pr2}The value on the right-hand side of \eqref{AA} on any basis vector from Group~${\rm II}$ equals zero.
\end{prop}
\begin{proof}Group ${\rm II}$ consists of vectors
\begin{gather*}
w=\frac{f}{T^{i_1}}\frac{f}{T^{i_2}}\frac{h}{T^{j_1}} \cdots \frac{h}{T^{j_s}} \frac{e}{T^l}v.
\end{gather*}

\begin{lem}\label{Group 2} We have
\begin{gather}
\bigg\langle \frac{f}{T^{a-1}} (v)^*, w \bigg\rangle = 2^{s+1} (m - lk), \label{G21} \\
\bigg\langle \frac{h}{T^{\ell}} \left(\frac{f}{T^{a-1-\ell}}v\right)^*, w \bigg\rangle =
 \begin{cases}
 2^{s+2} (m - lk), & \text{if $i_1 = i_2 = a-1-\ell$},\\
 2^{s+1} (m - lk), & \text{if $i_1 \ne i_2$ and $i_{1}$ or $i_2 = a-1-\ell$}, \\
 0, & \text{otherwise},
 \end{cases} \label{G22} \\
\bigg\langle \frac{e}{T^{\ell}}
\mathop{\sum_{i+j=a-1-\ell}}_{i\geq j\geq 0} \left(\frac{f}{T^i}\frac{f}{T^j}v\right)^*, w \bigg\rangle =
 \begin{cases}
 - 2^{s} (m - lk), & \text{if $\ell = a-1 - i_1 - i_2$},\\
 0, & \text{otherwise}.
 \end{cases} \label{G23}
\end{gather}
\end{lem}

\begin{proof}We have
\begin{gather}
\bigg\langle \frac{f}{T^{a-1}} (v)^*, w \bigg\rangle = \big\langle (v)^*, eT^{a-1} w \big\rangle = \bigg\langle (v)^*, \left[ hT^{a-1-i_1} +
\frac{f}{T^{i_1}} eT^{a-1} \right]
 \frac{f}{T^{i_2}}\frac{h}{T^{j_1}} \cdots \frac{h}{T^{j_s}} \frac{e}{T^l}v \bigg\rangle \nonumber \\
{}= \bigg\langle (v)^*, hT^{a-1-i_1}
\frac{f}{T^{i_2}}\frac{h}{T^{j_1}} \cdots \frac{h}{T^{j_s}} \frac{e}{T^l}v \bigg\rangle +
\bigg\langle (v)^*, \frac{f}{T^{i_1}} eT^{a-1}
 \frac{f}{T^{i_2}}\frac{h}{T^{j_1}} \cdots \frac{h}{T^{j_s}} \frac{e}{T^l}v \bigg\rangle. \label{he}
\end{gather}
Note that $eT^{a-1} \frac{f}{T^{i_2}} \frac{h}{T^{j_1}} \cdots \frac{h}{T^{j_s}} \frac{e}{T^l}$ is of degree $(-i_1 - 1, -i_1) \leq (-1,0)$, hence \newline $eT^{a-1} \frac{f}{T^{i_2}} \frac{h}{T^{j_1}}\cdots\frac{h}{T^{j_s}} \frac{e}{T^{l}} v = 0$. In the first summand in~(\ref{he}) we pull $hT^{a-1-i_1}$ to the right to get
\begin{gather*}
\bigg\langle (v)^*, -2 fT^{a-1-i_1-i_2} \frac{h}{T^{j_1}} \cdots \frac{h}{T^{j_s}} \frac{e}{T^l}v \bigg\rangle = \dots = \bigg\langle (v)^*, -2^{s+1} fT^{a-1-i_1-i_2-j_1-\dots -a_s} \frac{e}{T^l}v \bigg\rangle \\
\qquad{} = \bigg\langle (v)^*, -2^{s+1} fT^{l} \frac{e}{T^l}v \bigg\rangle = \big\langle (v)^*, 2^{s+1} (h - lc) v \big\rangle = 2^{s+1} (m - lk),
\end{gather*}
where at each step we do not write monomials of negative degree, since they give zero when applied to~$v$. This proves formula~(\ref{G21}).

We have
\begin{gather*}
\bigg\langle \frac{h}{T^\ell} \left(\frac{f}{T^{a-1-\ell}}v\right)^*, w \bigg\rangle = \bigg\langle \left(\frac{f}{T^{a-1-\ell}}v\right)^*, hT^{\ell} w \bigg\rangle,\\
hT^{\ell} w = \left[ -2fT^{\ell-i_1} + \frac{f}{T^{i_1}} hT^{\ell} \right] \frac{f}{T^{i_2}}\frac{h}{T^{j_1}} \cdots \frac{h}{T^{j_s}} \frac{e}{T^l} v \\
\hphantom{hT^{\ell} w} = -2fT^{\ell-i_1} \frac{f}{T^{i_2}}\frac{h}{T^{j_1}} \cdots \frac{h}{T^{j_s}} \frac{e}{T^l} v + \frac{f}{T^{i_1}} \left[ -2fT^{\ell-i_2} + \frac{f}{T^{i_2}} hT^{\ell} \right] \frac{h}{T^{j_1}} \cdots \frac{h}{T^{j_s}} \frac{e}{T^l} v \\
\hphantom{hT^{\ell} w} = \left[ -2 \frac{f}{T^{i_2}} fT^{\ell-i_1} - 2 \frac{f}{T^{i_1}}fT^{\ell-i_2} + \frac{f}{T^{i_1}} \frac{f}{T^{i_2}} hT^{\ell} \right] \frac{h}{T^{j_1}} \cdots \frac{h}{T^{j_s}} \frac{e}{T^l} v.
\end{gather*}
Note that the vector $\frac{f}{T^{i_1}} \frac{f}{T^{i_2}} hT^{\ell} \frac{h}{T^{j_1}} \cdots \frac{h}{T^{j_s}} \frac{e}{T^l} v$ after pulling $hT^{\ell}$ to the right either becomes a~zero vector or a~vector with two $f$'s, which of course gives zero when pairing with a basis vector with one $f$. Also, note that the only possibility for the vector $\frac{f}{T^{i_2}} fT^{\ell-i_1} \frac{h}{T^{j_1}} \cdots \frac{h}{T^{j_s}} \frac{e}{T^l} v$ to give a~nonzero pairing with $\big(\frac{f}{T^{a-1-\ell}}v\big)^*$ is when $i_2 = a-1-\ell$. Similarly $\frac{f}{T^{i_1}} fT^{\ell-i_2} \frac{h}{T^{j_1}} \cdots \frac{h}{T^{j_s}} \frac{e}{T^l} v$ gives a nonzero number only if $i_1 = a-1-\ell$. First consider the case $i_1=i_2=a-1-\ell$. We have
\begin{gather*}
- 4 \frac{f}{T^{a-1-\ell}} fT^{2\ell-a+1} \frac{h}{T^{j_1}} \cdots \frac{h}{T^{j_s}} \frac{e}{T^l} v \\
\qquad{} = - 4 \frac{f}{T^{a-1-\ell}} \left[ 2 fT^{2\ell-a+1-j_1} + \frac{h}{T^{j_1}} fT^{2\ell-a+1} \right] \frac{h}{T^{j_2}} \cdots \frac{h}{T^{j_s}} \frac{e}{T^l} v.
\end{gather*}
Note that $fT^{2\ell-a+1} \frac{h}{T^{j_2}} \cdots \frac{h}{T^{j_s}} \frac{e}{T^l}$ is of degree $(-j_1,-j_1) \leq (-1,-1)$, hence
\begin{gather*} fT^{2\ell-a+1} \frac{h}{T^{j_2}} \cdots \frac{h}{T^{j_s}} \frac{e}{T^l} v =0.\end{gather*} So we get
\begin{gather*}
 -8 \frac{f}{T^{a-1-\ell}} fT^{2\ell-a+1-j_1} \frac{h}{T^{j_2}} \cdots \frac{h}{T^{j_s}} \frac{e}{T^l} v \cdots
 = -2^{s+2} \frac{f}{T^{a-1-\ell}} fT^{2\ell-a+1-j_1-\cdots-j_s} \frac{e}{T^l} v \\
\qquad{} = - 2^{s+2} \frac{f}{T^{a-1-\ell}} fT^{l} \frac{e}{T^l} v = 2^{s+2} \frac{f}{T^{a-1-\ell}} (h - lc) v = 2^{s+2} (m - lk),
\end{gather*}
where at each step we don't write monomials of negative degree, since they give zero when applied to $v$.

For $i_1 = a-1-\ell \ne i_2$ and $i_2 = a-1-\ell \ne i_1$ we have
\begin{gather*}
- 2 \frac{f}{T^{a-1-\ell}} fT^{2\ell-a+1} \frac{h}{T^{j_1}} \dots \frac{h}{T^{j_s}} \frac{e}{T^l} v = 2^{s+1} (m - lk),
\end{gather*}
where we performed the exact same computation as above. Therefore, formula (\ref{G22}) is proved.

We have
\begin{gather*}
 \bigg\langle \frac{e}{T^{\ell}} \mathop{\sum_{i+j=a-1-\ell}}_{i\geq j\geq 0} \left(\frac{f}{T^i}\frac{f}{T^j}v\right)^*, w \bigg\rangle = \bigg\langle \mathop{\sum_{i+j=a-1-\ell}}_{i\geq j\geq 0} \left(\frac{f}{T^i}\frac{f}{T^j}v\right)^*, fT^{\ell} w \bigg\rangle, \\
fT^{\ell} w = \frac{f}{T^{i_1}}\frac{f}{T^{i_2}} fT^{\ell} \frac{h}{T^{j_1}} \cdots \frac{h}{T^{j_s}} \frac{e}{T^l} v = \frac{f}{T^{i_1}}\frac{f}{T^{i_2}} \left[ 2 fT^{\ell-j_1} +\frac{h}{T^{j_1}} fT^\ell \right] \frac{h}{T^{j_2}} \cdots \frac{h}{T^{j_s}} \frac{e}{T^l} v.
\end{gather*}
The only nonzero pairing happens when $\ell$ is such that $i_1+i_2 = a-1-\ell$. In that case $fT^{\ell} \frac{h}{T^{j_2}} \cdots \frac{h}{T^{j_s}} \frac{e}{T^l}$ has degree $(-j_1, -j_1) \leq (-1,-1)$, hence $fT^{\ell} \frac{h}{T^{j_2}} \cdots \frac{h}{T^{j_s}} \frac{e}{T^l}v=0$. Therefore we have
\begin{gather*}
2 \frac{f}{T^{i_1}}\frac{f}{T^{i_2}} fT^{\ell-j_1} \frac{h}{T^{j_2}} \cdots \frac{h}{T^{j_s}} \frac{e}{T^l} v = 2^s \frac{f}{T^{i_1}}\frac{f}{T^{i_2}} fT^{\ell-j_1 - \dots - j_s} \frac{e}{T^l} v = 2^s \frac{f}{T^{i_1}}\frac{f}{T^{i_2}} fT^{l} \frac{e}{T^l} v,
\end{gather*}
where we pulled $fT^{\ell-j_1}$ to the right and did not write monomials of negative degree, since they give zero when applied to $v$. Hence we get
\begin{gather*}
2^s \frac{f}{T^{i_1}}\frac{f}{T^{i_2}} (-h +lc) v.
\end{gather*}
Therefore,
\begin{gather*}
\bigg\langle \frac{e}{T^{\ell}} \mathop{\sum_{i+j=a-1-\ell}}_{i\geq j\geq 0} \left(\frac{f}{T^i}\frac{f}{T^j}v\right)^*, w \bigg\rangle \\
\qquad{} = \mathop{\sum_{i+j=a-1-\ell}}_{i\geq j\geq 0}\bigg\langle \left(\frac{f}{T^i}\frac{f}{T^j}v\right)^*, 2^s \frac{f}{T^{i_1}}\frac{f}{T^{i_2}} (-h +lc) v \bigg\rangle = - 2^{s} (m - lk),
\end{gather*}
since for $i = i_1$, $j = i_2$ we get $-2^s (m - lk)$ and zero for other pairs $(i,j)$. Formula (\ref{G23}) and Lemma~\ref{Group 2} are proved.
\end{proof}

By Lemma \ref{Group 2}, we have
\begin{gather*}
\bigg\langle \frac{f}{T^{a-1}} (v)^*-
\sum_{\ell=1}^{a-1}\bigg[\frac{h}{T^{\ell}}\left(
\frac{f}{T^{a-1-\ell}}v\right)^* + 2\frac{e}{T^{\ell}} \mathop{\sum_{i+j=a-1-\ell}}_{i\geq j\geq 0} \left(\frac{f}{T^i}\frac{f}{T^j}v\right)^*\bigg], w \bigg\rangle \\
\qquad{} = 2^{s+1} (m-lk) - 2 \cdot 2^{s+1} (m-lk) + 2 \cdot 2^{s} (m-lk) = 0.
\end{gather*}
Note that
\begin{gather*}
 \bigg\langle \sum_{\ell=1}^{a-1} \frac{h}{T^{\ell}}\left(\frac{f}{T^{a-1-\ell}}v\right)^*, w \bigg\rangle = 2 \cdot 2^{s+1} (m - lk)
\end{gather*}
in both cases $i_1 = i_2$ and $i_1 \ne i_2$. Also note that
\begin{gather*}
\bigg\langle \sum_{\ell=1}^{a-1}\ \frac{e}{T^{\ell}} \mathop{\sum_{i+j=a-1-\ell}}_{i\geq j\geq 0} \left(\frac{f}{T^i}\frac{f}{T^j}v\right)^*, w \bigg\rangle \ne 0
\end{gather*}
only if $\ell$ is such that $i_1+i_2 = a-1-\ell$. Therefore, Proposition \ref{Pr2} is proved.
\end{proof}

\subsection[Group ${\rm III}$]{Group $\boldsymbol{{\rm III}}$}
\begin{prop}\label{Pr1}The value of the right-hand side of \eqref{AA} on any basis vector of Group~${\rm III}$ equals zero.
\end{prop}
\begin{proof}A vector in Group ${\rm III}$ has the form
\begin{gather*}
w= \frac{f}{T^{i_1}}\cdots\frac{f}{T^{i_r}}
\frac{h}{T^{j_1}}\cdots\frac{h}{T^{j_s}}
\frac{e}{T^{l_1}}\cdots\frac{e}{T^{l_{r-1}}}v,
\end{gather*}
where $r \geq 3$.

\begin{lem}\label{Group 1}For every $\ell \in \{1, \dots, a-1 \}$, we have
\begin{gather}
\bigg\langle \frac{f}{T^{a-1}} (v)^*, w \bigg\rangle =0, \label{G31} \\
\bigg\langle \frac{h}{T^{\ell}} \left(\frac{f}{T^{a-1-\ell}}v\right)^*, w \bigg\rangle =0, \label{G32} \\
 \bigg\langle \frac{e}{T^{\ell}} \mathop{\sum_{i+j=a-1-\ell}}_{i\geq j\geq 0} \left(\frac{f}{T^i}\frac{f}{T^j}v\right)^*, w \bigg\rangle = 0. \label{G33}
\end{gather}
\end{lem}

\begin{proof}We have
\begin{gather*}
 \bigg\langle \frac{f}{T^{a-1}} (v)^*, w \bigg\rangle = \big\langle (v)^*, eT^{a-1} w \big\rangle \\
\hphantom{\bigg\langle \frac{f}{T^{a-1}} (v)^*, w \bigg\rangle} = \bigg\langle (v)^*, \left[hT^{a-1-i_1} +\frac{f}{T^{i_1}} eT^{a-1}\right]
 \frac{f}{T^{i_2}} \cdots \frac{f}{T^{i_r}} \frac{h}{T^{j_1}}\cdots\frac{h}{T^{j_s}} \frac{e}{T^{l_1}}\cdots\frac{e}{T^{l_{r-1}}} v \bigg\rangle.
\end{gather*}
Note that $eT^{a-1} \frac{f}{T^{i_2}} \cdots \frac{f}{T^{i_r}} \frac{h}{T^{j_1}}\cdots\frac{h}{T^{j_s}} \frac{e}{T^{l_1}}\cdots\frac{e}{T^{l_{r-1}}} v $ is of degree $(-i_1 - 1,-i_1) \leq (-1,0)$, hence zero. So we have
\begin{gather*}
\bigg\langle (v)^*, \left[-2 fT^{a-1-i_1-i_2} +\frac{f}{T^{i_2}} hT^{a-1-i_1} \right]
\frac{f}{T^{i_3}} \cdots \frac{f}{T^{i_r}} \frac{h}{T^{j_1}}\cdots
\frac{h}{T^{j_s}} \frac{e}{T^{l_1}}\cdots\frac{e}{T^{l_{r-1}}} v \bigg\rangle.
\end{gather*}
As above, note that $hT^{a-1-i_1} \frac{f}{T^{i_3}} \cdots \frac{f}{T^{i_r}} \frac{h}{T^{j_1}}\cdots\frac{h}{T^{j_s}} \frac{e}{T^{l_1}}\cdots\frac{e}{T^{l_{r-1}}} v $ is of degree $ (-i_2-1, -i_2) \leq (-1,0)$, hence zero. Therefore we obtain
\begin{gather*}
\bigg\langle (v)^*, -2 \frac{f}{T^{i_3}} \cdots \frac{f}{T^{i_r}} fT^{a-1-i_1-i_2} \frac{h}{T^{j_1}}\cdots\frac{h}{T^{j_s}} \frac{e}{T^{l_1}}\cdots\frac{e}{T^{l_{r-1}}} v \bigg\rangle = 0,
\end{gather*}
since $fT^{a-1-i_1-i_2} \frac{h}{T^{j_1}}\cdots\frac{h}{T^{j_s}} \frac{e}{T^{l_1}}\cdots\frac{e}{T^{l_{r-1}}} v $ is of degree $(-i_3 - \dots - i_r - r + 2, -i_2 - \dots - i_r) \leq (-1,0)$ for $r \geq 3$, hence zero. Formula (\ref{G31}) is proved.

We have
\begin{gather}
 \bigg\langle \frac{h}{T^\ell} \left(\frac{f}{T^{a-1-\ell}}v\right)^*, w \bigg\rangle = \bigg\langle \left(\frac{f}{T^{a-1-\ell}}v\right)^*, hT^{\ell} w \bigg\rangle,\nonumber\\
hT^{\ell} w = \left[-2 fT^{\ell-i_1}+ \frac{f}{T^{i_1}} hT^{\ell} \right] \frac{f}{T^{i_2}}\cdots\frac{f}{T^{i_r}} \frac{h}{T^{j_1}}\cdots\frac{h}{T^{j_s}} \frac{e}{T^{l_1}}\cdots\frac{e}{T^{l_{r-1}}} v \nonumber \\
\hphantom{hT^{\ell} w} = -2 \frac{f}{T^{i_2}}\cdots\frac{f}{T^{i_r}} fT^{\ell-i_1} \frac{h}{T^{j_1}}\cdots\frac{h}{T^{j_s}} \frac{e}{T^{l_1}}\cdots\frac{e}{T^{l_{r-1}}} v \label{fr1} \\
\hphantom{hT^{\ell} w=} + \frac{f}{T^{i_1}} hT^{\ell} \frac{f}{T^{i_2}}\cdots\frac{f}{T^{i_r}} \frac{h}{T^{j_1}}\cdots\frac{h}{T^{j_s}} \frac{e}{T^{l_1}}\cdots\frac{e}{T^{l_{r-1}}} v. \label{fr2}
\end{gather}
Observe that for $\ell \leq i_1$ in~(\ref{fr1}) we have a basis vector, hence it gives zero when pairing with $\big(\frac{f}{T^{a-1-\ell}}v\big)^*$. If $\ell > i_1$, then we pull $fT^{\ell-i_1}$ to the right and notice that no matter how $fT^{\ell-i_1}$ interacts with $h$'s and $e$'s, it does not affect the number of~$f$'s, which is greater or equal than two. Hence, the vector in~(\ref{fr1}) gives zero when pairing with $\big(\frac{f}{T^{a-1-\ell}}v\big)^*$.

In (\ref{fr2}) note that
\begin{gather*}
hT^{\ell} \frac{f}{T^{i_2}} = -2 fT^{\ell-i_2} + \frac{f}{T^{i_2}} hT^{\ell}
\end{gather*}
so that either $hT^{\ell}$ is pulled to the right not affecting the number of $f$'s or it gives $fT^{\ell-i_2}$, for which we apply the same argument as above after pulling it to the right to argue that the pairing of the vector in~(\ref{fr2}) with $\big(\frac{f}{T^{a-1-\ell}}v\big)^*$ is zero. Formula (\ref{G32}) is proved.

We have
\begin{gather*}
\bigg\langle \frac{e}{T^{\ell}} \mathop{\sum_{i+j=a-1-\ell}}_{i\geq j\geq 0} \left(\frac{f}{T^i}\frac{f}{T^j}v\right)^*, w \bigg\rangle = \bigg\langle \mathop{\sum_{i+j=a-1-\ell}}_{i\geq j\geq 0} \left(\frac{f}{T^i}\frac{f}{T^j}v\right)^*, fT^{\ell} w \bigg\rangle, \\
fT^{\ell} w = fT^{\ell} \frac{f}{T^{i_1}}\cdots\frac{f}{T^{i_r}}
\frac{h}{T^{j_1}}\cdots\frac{h}{T^{j_s}} \frac{e}{T^{l_1}}\cdots\frac{e}{T^{l_{r-1}}} v \\
\hphantom{fT^{\ell} w} = \frac{f}{T^{i_1}}\cdots\frac{f}{T^{i_r}} fT^{\ell}
\frac{h}{T^{j_1}}\cdots\frac{h}{T^{j_s}} \frac{e}{T^{l_1}}\cdots\frac{e}{T^{l_{r-1}}} v.
\end{gather*}
As in formula (\ref{G32}), no matter how $fT^{\ell}$ interacts with $h$'s and $e$'s, the number of $f$'s remains unchanged, i.e., we have more than or equal to three~$f$'s, so that pairing with $\big(\frac{f}{T^i}\frac{f}{T^j}v\big)^*$ is zero. Formula~(\ref{G33}) is proved.
\end{proof}

Proposition \ref{Pr1} follows from Lemma \ref{Group 1}.
\end{proof}

Theorem \ref{thm form} is proved.
